\is{Ablautklasse| see {Verbflexion, Ablautklasse}}
\is{Abtönungspartikel| see {Partikel!Abtönungspartikel}}
\is{Akkumulierbarkeitstest|see {Ergänzung, Akkumulierbarkeitstest}}
\is{Aktiv| see {Genus verbi, Aktiv}}
\is{Anhebung| see {Verbkomplex, Anhebung}}
\is{Assoziationsstärke|see {Ergänzung, Assoziationsstärke}}
\is{atelisch|see {Verbklassen, atelisch}}
\is{Ausklammerung| see {Nachfeld, Ausklammerung}}
\is{Auslagerungstest|see {Ergänzung, Auslagerungstest}}
\is{Aussagesatz| seealso {Verbzweitsatz}}
\is{autonome Rollenkodierung|see {Angabe, autonome Rollenkodierung}}
\issa{Basisgraphem}{Graphem}
\is{Betonungszeichen| see{Transkription, Betonungszeichen}}
\is{ch-Laute| see{Laute, ch-Laute}}
\is{CVC| see{Silbe, CVC}}
\is{CV| see{Silbe, CV}}
\is{Dativus Iudicantis| seealso {Kommentaradverbial, Dativus Iudicantis}}
\is{Dependens| see {Dependenz, Dependens}} 
\is{Derivationsstammform| see {Stammform!Derivation}}
\is{Determinans| see {Determinativkompositum!Determinans}}
\is{Determinatum| see {Determinativkompositum!Determinatum}}
\is{direktes Objekt| see {Akkusativobjekt}}
\is{Eigenname| see {Substantiv, Eigenname}}
\is{einfach| see{Graphem, einfach}}
\is{Ellipse| see {Koordination!Ellipse}}
\is{eng| see{Transkription, eng}}
\is{Entscheidungsfragesatz| seealso {Verberstsatz}}
\is{Erfragbarkeitstest| see {Satzkonstituente, Erfragbarkeitstest}}
\is{Ergänzungsfragesatz| seealso {Verbzweitsatz}}
\is{Ersetzungstest| see {Satzkonstituente, Ersetzungstest}}
\is{fakulativ| see{phonologischer Prozess, fakulatitiv}}
\is{Flexionsstammform| see {Stammform!lexikalisch}}
\is{Fokuspartikel| see {Partikel!Fokuspartikel}}
\is{freier Dativ| see {Dativobjekt, dynamisch}}
\is{frei| see{Allophon, frei}}
\is{Futur II| see {Tempus, Futur II}}
\is{Futur I| see {Tempus, Futur I}}
\is{Gattungsbezeichnung| see {Substantiv, Gattungsbezeichnung}}
\is{geschlossen| see{Silbe, geschlossen}}
\is{graphematisch komplex| see{Onset, graphematisch komplex}}
\is{Halbvokal| see{Approximant}}
\is{Handlungsverb|see {Verbklassen, Handlungsverb}}
\is{harter Gaumen| see{Palatum}}
\is{Hauptakzent| see{Akzent, Hauptakzent}}
\is{h@[h]| see{Laute, [h]}}
\is{Imperativsatz| seealso {Verberstsatz}}
\is{Imperativ| see {Modus, Imperativ}}
\is{Implikationstest|see {Ergänzung, Implikationstest}}
\is{Indefinitpronomen| see {Pronomen, Indefinitpronomen}}
\is{Indikativ| see {Modus, Indikativ}}
\is{indirektes Objekt| see {Dativobjekt}}
\is{inkohärent konstruierend| see {Verbkomplex, inkohärent konstruierend}}
\is{Interrogativpronomen| see {Pronomen, \textit{w}-Pronomen}}
\is{Klammern| see{Transkription, Klammern}}
\is{Knacklaut| see{Laute, Knacklaut}}
\is{kognates Objekt| see {Akkusativobjekt, dynamisch}}
\is{kohärent konstruierend| see {Verbkomplex, kohärent konstruierend}}
\is{Komparation| see {Adjektiv, Komparation}}
\is{komplementär| see{Allophon, komplementär}}
\is{komplex| see{Graphem, komplex}}
\is{Kompositionsstammform| see {Stammform!Komposition}}
\is{Konjunktionaladverb| see {Adverb, kausal}}
\is{Konjunktiv II| see {Modus, Konjunktiv II}}
\is{Konjunktiv I| see {Modus, Konjunktiv I}}
\is{Kontrolle| see {Verbkomplex, Kontrolle}}
\is{Koordinationstest| see {Satzkonstituente, Koordinationstest}}
\is{Kopf| seealso {Phrase}}
\is{Längenzeichen| see{Transkription, Längenzeichen}}
\is{lateral| see{Approximant, lateral}}
\is{Listem| see {Wortbegriff, Listem}}
\is{Mitlaut| seealso{Konsonant}}
\is{nasal| see{Vokal, nasal}}
\is{Nebenakzent| see{Akzent, Nebenakzent}}
\is{Nebensatz| seealso {Verbletztsatz}}
\is{Oberfeld| see {Verbletztsatz, Oberfeld}}
\is{obligatorisch| see{phonologischer Prozess, obligatorisch}}
\is{offen| see{Silbe, offen}}
\issa{Orthographem}{Graphem}
\is{Passiv| seealso {Genus verbi, Passiv}}
\is{Perfekt| see {Tempus, Perfekt}}
\is{Personalpronomen| see {Pronomen, Personalpronomen}}
\is{Platzhalter| see {\textit{es}, Platzhalter}}
\is{Plusquamperfekt| see {Tempus, Plusquamperfekt}}
\is{Possessivpronomen| see {Pronomen, Possessivpronomen}}
\is{Präsens| see {Tempus, Präsens}}
\is{Präteritum| see {Tempus, Präteritum}}
\is{progressiv| see{Assimilation, progressiv}}
\is{Pronominaladverb| see {Adverb, Präpositionaladverb}}
\is{Pronominalisierungstest| see {Satzkonstituente, Pronominalisierungstest}}
\is{r-Diphthong| see{Diphthong, zentralisierend}}
\is{Reflexivpronomen| seealso {Pronomen, reflexiver Gebrauch}}
\is{Regens| see {Dependenz, Regens}} 
\is{regressiv| see{Assimilation, regressiv}}
\is{Relativpronomen| see {Pronomen, im Relativsatz}}
\is{Restriktionen|see {Angabe, Restriktionen}}
\is{Rollenzuweisung|see {Ergänzung, Rollenzuweisung}}
\is{Satzäquivalent| see {Partikel!Satzäquivalent}}
\is{Satzglied| seealso {Satzkonstituente}}
\is{Satzklammer| seealso {Verbklammer}}
\is{schließend| see{Diphthong, schließend}}
\is{Schrägstriche| see{Transkription, Schrägstriche}}
\is{Schwa| see{Laute, Schwa}}
\is{Selbstlaut| see{Vokal}}
\is{Silbenanlautgesetz| see{Silbenpräferenzgesetz, Silbenanlautgesetz}}
\is{Silbenauslautgesetz| see{Silbenpräferenzgesetz, Silbenauslautgesetz}}
\is{Silbenkerngesetz| see{Silbenpräferenzgesetz, Silbenkerngesetz}}
\is{silbischer Konsonant| see{Transkription, silbischer Konsonant}}
%\is{some term| seealso {some other term}}
%\is{some term| see {some other term}}
\is{speziell| see{Silbengelenksschreibung, speziell}}
\is{s@[s]| seealso{Laute, \textipa{[z]}}}
\is{s@[s]| see{Laute, [s]}}
\is{Steigerungspartikel| see {Partikel!Steigerungspartikel}}
\is{stellungsbedingt| see{Allophon, stellungsbedingt}}
\is{Subjunktionalsatz| see {Nebensatz, Subjunktionalsatz}}
\is{Subkategorisierungstest|see {Ergänzung, Subkategorisierungstest}}
\is{Substantivgroßschreibung| see{satzinterne Großschreibung}}
\is{Tätigkeitsverb|see {Verbklassen, Tätigkeitsverb}}
\is{telisch|see {Verbklassen, telisch}}
\is{dz@\textipa{[DZ]}| see{Laute, \textipa{[DZ]}}}
\is{z@\textipa{[Z]}| see{Laute, \textipa{[Z]}}}
\is{tiefes Schwa| see{Laute, tiefes Schwa}}
\is{transitiv|see {Verbklassen, transitiv}}
\is{Unterfeld| see {Verbletztsatz, Unterfeld}}
\is{Unterspezifiziertheit| see {Determinativkompositum!Unterspezifiziertheit}}
\is{Verbletztsatz| seealso {Nebensatz}}
\is{Voranstellungstest| see {Satzkonstituente, Voranstellungstest}}
\is{Vorgangsverb|see {Verbklassen, Vorgangsverb}}
\is{Wechselpräposition| see {Präposition!Bewegung vs. Position}}
\is{Weglassbarkeitstest|see {Ergänzung, Weglassbarkeitstest}}
\is{weicher Gaumen| see{Velum}}
\is{weit| see{Transkription, weit}}
\is{Wortakzent| see{Akzent, Wortakzent}}
\is{Zählbarkeit| see {Substantiv, Zählbarkeit}}
\is{Zahndamm| see{Alveolen}}
\is{Zäpfchen| see{Uvula}}
\is{zentralisierend| see{Diphthong, zentralisierend}}
\is{zentral| see{Approximant, zentral}}
\is{Zustandsverb|see {Verbklassen, Zustandsverb}}
\is{Zwielaut| see{Diphthong}} %!!! nachsehen, ob es im Kapitel falsch ist
